\documentclass[pdflatex,sn-nature]{sn-jnl}
\usepackage{amsmath,amssymb,amsfonts}
\usepackage{booktabs}
\usepackage{graphicx}
\begin{document}
\title{Supplementary Information for TSRClust}
\author{Hasan Bulut}
\maketitle
\section*{Supplementary Methods}
The reproducibility files contain the complete simulation manifests and replication-level result exports. The architecture study used 80 cells and 35,200 generated datasets; the stress study used 30 cells and 9,000 generated datasets; and the expanded negative-control study used 10 cells and 10,000 generated datasets.
\section*{S1. Complete simulation design}
Architecture, stress, and negative-control cell-summary CSV files provide compact cell-level summaries. Full replication-level results are provided with the reproducibility package rather than reproduced as large tables here.
\section*{S2. Paired comparisons}
Architecture, stress, and negative-control paired-comparison CSV files contain method differences against TSRClust with Monte Carlo standard errors.
\begin{table}[!htbp]
\caption{Validation ablation summary for architecture and stress simulations.}
\label{tab:supp-validation-ablation}
\centering
\tiny
\setlength{\tabcolsep}{1pt}
\begin{tabular}{llrrrrr}
\hline
Study & Method & MajorARI & Min. F1 & Any & AbsQ & Spur. \\
\hline
Architecture & TSRClust & 0.946 & 0.715 & 0.877 & 0.280 & 0.117 \\
Architecture & TSRClust w/o val. & 0.946 & 0.606 & 1.000 & 0.567 & 0.495 \\
Stress & TSRClust & 0.720 & 0.535 & 0.639 & 0.448 & 0.014 \\
Stress & TSRClust w/o val. & 0.720 & 0.461 & 1.000 & 0.127 & 0.312 \\
\hline
\end{tabular}
\begin{flushleft}\footnotesize Here, w/o val. denotes without validation. The architecture ablation summary pools 19,200 positive-scenario simulation replications (M1--M4), while the stress summary is based on 9,000 simulation replications. Detection F1 and runtime are included in the companion CSV file. The ablation removes only the Monte Carlo validation layer while retaining the same Stage-I residual set and Stage-II DBSCAN scan. Numerical values are copied from the final summary result files; no simulations are rerun.\end{flushleft}
\end{table}

\section*{S3. Monte Carlo uncertainty}
For expanded negative controls, TSRClust produced 3 false micro-discoveries in 10000 successful evaluations, giving an estimated probability of 0.0003 with an exact 95\% binomial confidence interval of [0.00006, 0.00088]. TSRClust without validation produced a false micro-discovery in every negative-control run.
\section*{S4. OTRIMLE sensitivity and runtime}
OTRIMLE sensitivity outputs and detailed runtime summaries are included in the reproducibility results folders. The main text reports the fixed-\(K\) OTRIMLE results used in the headline comparison.
\section*{S5. Real-data details}
Real-data supplementary CSV files provide class-specific trimming rates and accepted residual micro-cluster compositions for the olive-oil and Shuttle analyses.
\begin{figure}[htbp]
\centering
\includegraphics[width=\linewidth]{../figures/fig6_realdata_microclusters.pdf}
\caption{Accepted real-data TSRClust micro-clusters. Bar heights show accepted residual micro-cluster sizes; fill identifies the dominant external class.}
\label{fig:supp-realdata-microclusters}
\end{figure}
\end{document}
